%% kapa: exam=1 round=19 year=2008 period=1 form=A subject=appraisal_law subject_ko=부동산관계법규 seq=14 number=54
%% source: Q-Net 감정평가사 기출문제 (cst00309) artlSeq=1001962 file=extracted/1008/1001962/2008년 제19회 감정평가사 1차시험 기출문제/1교시.hwp
%% conversion: kordoc 4.24.51 -> tools/gen1.py (auto); review=proofread-sonnet; flags=none
\begin{question}{54}
국토의 계획 및 이용에 관한 법령상 토지거래허가에 관한 설명으로 옳은 것은?
\begin{choices}
  \item 토지거래허가구역을 지정하기 위해서는 관계 행정기관의 장과 반드시 협의를 하여야 하고, 또한 도시계획위원회의 심의를 거쳐야 한다.
  \item 허가대상 토지거래계약이라 함은 허가구역 안 또는 일정한 경우 인접지역에 있는 토지의 소유권·지상권을 이전 또는 설정하는 계약행위를 말한다.
  \item 법령의 제정·개정 또는 폐지나 그에 의한 고시·공고로 인하여 토지이용에 대한 행위제한이 강화되거나 해제되는 지역은 토지거래허가구역의 지정대상지역에 속한다.
  \item 중심상업지역에서 허가구역 지정 당시 800m\textsuperscript{2}인 토지를 토지거래허가구역의 지정 후 각각 200m\textsuperscript{2}의 4필지로 분할한 경우에는 분할 이후 당해 분할된 토지에 대한 최초의 거래는 허가를 받아야 한다.
  \item 국토해양부장관이 허가구역을 지정한 때에는 지체없이 허가구역의 범위, 지정기간, 허가를 요하는 토지의 면적을 공고하고, 그 공고내용을 시·도지사에게 통지하여야 하며, 허가구역의 지정 및 재지정, 해제, 축소의 경우의 각각의 효력은 공고한 날부터 5일이 경과한 후에 발생한다.
\end{choices}
\end{question}
\answer{4}
